\IfFileExists{stacks-project.cls}{%
\documentclass{stacks-project}
}{%
\documentclass{amsart}
}

% For dealing with references we use the comment environment
\usepackage{verbatim}
\newenvironment{reference}{\comment}{\endcomment}
%\newenvironment{reference}{}{}
\newenvironment{slogan}{\comment}{\endcomment}
\newenvironment{history}{\comment}{\endcomment}

% For commutative diagrams we use Xy-pic
\usepackage[all]{xy}

% We use 2cell for 2-commutative diagrams.
\xyoption{2cell}
\UseAllTwocells

% We use multicol for the list of chapters between chapters
\usepackage{multicol}

% This is generally recommended for better output
\usepackage{lmodern}
\usepackage[T1]{fontenc}

% For cross-file-references
\usepackage{xr-hyper}

% Package for hypertext links:
\usepackage{hyperref}

% For any local file, say "hello.tex" you want to link to please
% use \externaldocument[hello-]{hello}
\externaldocument[introduction-]{introduction}
\externaldocument[conventions-]{conventions}
\externaldocument[sets-]{sets}
\externaldocument[categories-]{categories}
\externaldocument[topology-]{topology}
\externaldocument[sheaves-]{sheaves}
\externaldocument[sites-]{sites}
\externaldocument[stacks-]{stacks}
\externaldocument[fields-]{fields}
\externaldocument[algebra-]{algebra}
\externaldocument[brauer-]{brauer}
\externaldocument[homology-]{homology}
\externaldocument[derived-]{derived}
\externaldocument[simplicial-]{simplicial}
\externaldocument[more-algebra-]{more-algebra}
\externaldocument[smoothing-]{smoothing}
\externaldocument[modules-]{modules}
\externaldocument[sites-modules-]{sites-modules}
\externaldocument[injectives-]{injectives}
\externaldocument[cohomology-]{cohomology}
\externaldocument[sites-cohomology-]{sites-cohomology}
\externaldocument[dga-]{dga}
\externaldocument[dpa-]{dpa}
\externaldocument[sdga-]{sdga}
\externaldocument[hypercovering-]{hypercovering}
\externaldocument[schemes-]{schemes}
\externaldocument[constructions-]{constructions}
\externaldocument[properties-]{properties}
\externaldocument[morphisms-]{morphisms}
\externaldocument[coherent-]{coherent}
\externaldocument[divisors-]{divisors}
\externaldocument[limits-]{limits}
\externaldocument[varieties-]{varieties}
\externaldocument[topologies-]{topologies}
\externaldocument[descent-]{descent}
\externaldocument[perfect-]{perfect}
\externaldocument[more-morphisms-]{more-morphisms}
\externaldocument[flat-]{flat}
\externaldocument[groupoids-]{groupoids}
\externaldocument[more-groupoids-]{more-groupoids}
\externaldocument[etale-]{etale}
\externaldocument[chow-]{chow}
\externaldocument[intersection-]{intersection}
\externaldocument[pic-]{pic}
\externaldocument[weil-]{weil}
\externaldocument[adequate-]{adequate}
\externaldocument[dualizing-]{dualizing}
\externaldocument[duality-]{duality}
\externaldocument[discriminant-]{discriminant}
\externaldocument[derham-]{derham}
\externaldocument[local-cohomology-]{local-cohomology}
\externaldocument[algebraization-]{algebraization}
\externaldocument[curves-]{curves}
\externaldocument[resolve-]{resolve}
\externaldocument[models-]{models}
\externaldocument[functors-]{functors}
\externaldocument[equiv-]{equiv}
\externaldocument[pione-]{pione}
\externaldocument[etale-cohomology-]{etale-cohomology}
\externaldocument[proetale-]{proetale}
\externaldocument[relative-cycles-]{relative-cycles}
\externaldocument[more-etale-]{more-etale}
\externaldocument[trace-]{trace}
\externaldocument[crystalline-]{crystalline}
\externaldocument[spaces-]{spaces}
\externaldocument[spaces-properties-]{spaces-properties}
\externaldocument[spaces-morphisms-]{spaces-morphisms}
\externaldocument[decent-spaces-]{decent-spaces}
\externaldocument[spaces-cohomology-]{spaces-cohomology}
\externaldocument[spaces-limits-]{spaces-limits}
\externaldocument[spaces-divisors-]{spaces-divisors}
\externaldocument[spaces-over-fields-]{spaces-over-fields}
\externaldocument[spaces-topologies-]{spaces-topologies}
\externaldocument[spaces-descent-]{spaces-descent}
\externaldocument[spaces-perfect-]{spaces-perfect}
\externaldocument[spaces-more-morphisms-]{spaces-more-morphisms}
\externaldocument[spaces-flat-]{spaces-flat}
\externaldocument[spaces-groupoids-]{spaces-groupoids}
\externaldocument[spaces-more-groupoids-]{spaces-more-groupoids}
\externaldocument[bootstrap-]{bootstrap}
\externaldocument[spaces-pushouts-]{spaces-pushouts}
\externaldocument[spaces-chow-]{spaces-chow}
\externaldocument[groupoids-quotients-]{groupoids-quotients}
\externaldocument[spaces-more-cohomology-]{spaces-more-cohomology}
\externaldocument[spaces-simplicial-]{spaces-simplicial}
\externaldocument[spaces-duality-]{spaces-duality}
\externaldocument[formal-spaces-]{formal-spaces}
\externaldocument[restricted-]{restricted}
\externaldocument[spaces-resolve-]{spaces-resolve}
\externaldocument[formal-defos-]{formal-defos}
\externaldocument[defos-]{defos}
\externaldocument[cotangent-]{cotangent}
\externaldocument[examples-defos-]{examples-defos}
\externaldocument[algebraic-]{algebraic}
\externaldocument[examples-stacks-]{examples-stacks}
\externaldocument[stacks-sheaves-]{stacks-sheaves}
\externaldocument[criteria-]{criteria}
\externaldocument[artin-]{artin}
\externaldocument[quot-]{quot}
\externaldocument[stacks-properties-]{stacks-properties}
\externaldocument[stacks-morphisms-]{stacks-morphisms}
\externaldocument[stacks-limits-]{stacks-limits}
\externaldocument[stacks-cohomology-]{stacks-cohomology}
\externaldocument[stacks-perfect-]{stacks-perfect}
\externaldocument[stacks-introduction-]{stacks-introduction}
\externaldocument[stacks-more-morphisms-]{stacks-more-morphisms}
\externaldocument[stacks-geometry-]{stacks-geometry}
\externaldocument[moduli-]{moduli}
\externaldocument[moduli-curves-]{moduli-curves}
\externaldocument[examples-]{examples}
\externaldocument[exercises-]{exercises}
\externaldocument[guide-]{guide}
\externaldocument[desirables-]{desirables}
\externaldocument[coding-]{coding}
\externaldocument[obsolete-]{obsolete}
\externaldocument[fdl-]{fdl}
\externaldocument[index-]{index}

% Theorem environments.
%
\theoremstyle{plain}
\newtheorem{theorem}[subsection]{Theorem}
\newtheorem{proposition}[subsection]{Proposition}
\newtheorem{lemma}[subsection]{Lemma}

\theoremstyle{definition}
\newtheorem{definition}[subsection]{Definition}
\newtheorem{example}[subsection]{Example}
\newtheorem{exercise}[subsection]{Exercise}
\newtheorem{situation}[subsection]{Situation}

\theoremstyle{remark}
\newtheorem{remark}[subsection]{Remark}
\newtheorem{remarks}[subsection]{Remarks}

\numberwithin{equation}{subsection}

% Macros
%
\def\lim{\mathop{\mathrm{lim}}\nolimits}
\def\colim{\mathop{\mathrm{colim}}\nolimits}
\def\Spec{\mathop{\mathrm{Spec}}}
\def\Hom{\mathop{\mathrm{Hom}}\nolimits}
\def\Ext{\mathop{\mathrm{Ext}}\nolimits}
\def\SheafHom{\mathop{\mathcal{H}\!\mathit{om}}\nolimits}
\def\SheafExt{\mathop{\mathcal{E}\!\mathit{xt}}\nolimits}
\def\Sch{\mathit{Sch}}
\def\Mor{\mathop{\mathrm{Mor}}\nolimits}
\def\Ob{\mathop{\mathrm{Ob}}\nolimits}
\def\Sh{\mathop{\mathit{Sh}}\nolimits}
\def\NL{\mathop{N\!L}\nolimits}
\def\CH{\mathop{\mathrm{CH}}\nolimits}
\def\proetale{{pro\text{-}\acute{e}tale}}
\def\etale{{\acute{e}tale}}
\def\QCoh{\mathit{QCoh}}
\def\Ker{\mathop{\mathrm{Ker}}}
\def\Im{\mathop{\mathrm{Im}}}
\def\Coker{\mathop{\mathrm{Coker}}}
\def\Coim{\mathop{\mathrm{Coim}}}

% Boxtimes
%
\DeclareMathSymbol{\boxtimes}{\mathbin}{AMSa}{"02}

%
% Macros for moduli stacks/spaces
%
\def\QCohstack{\mathcal{QC}\!\mathit{oh}}
\def\Cohstack{\mathcal{C}\!\mathit{oh}}
\def\Spacesstack{\mathcal{S}\!\mathit{paces}}
\def\Quotfunctor{\mathrm{Quot}}
\def\Hilbfunctor{\mathrm{Hilb}}
\def\Curvesstack{\mathcal{C}\!\mathit{urves}}
\def\Polarizedstack{\mathcal{P}\!\mathit{olarized}}
\def\Complexesstack{\mathcal{C}\!\mathit{omplexes}}
% \Pic is the operator that assigns to X its picard group, usage \Pic(X)
% \Picardstack_{X/B} denotes the Picard stack of X over B
% \Picardfunctor_{X/B} denotes the Picard functor of X over B
\def\Pic{\mathop{\mathrm{Pic}}\nolimits}
\def\Picardstack{\mathcal{P}\!\mathit{ic}}
\def\Picardfunctor{\mathrm{Pic}}
\def\Deformationcategory{\mathcal{D}\!\mathit{ef}}
\usepackage{cleveref}
\crefname{originalequation}{original equation}{original equations}
\Crefname{originalequation}{Original equation}{Original equations}
\makeatletter
\newlabel{equation-compare}{{06UR}{source}{}{sourcepointer-1}{}}
\newlabel{equation-compare@cref}{{[originalequation][0][]06UR}{[1][0][]source}{}{}{}}
\makeatother
\title{The bounded-above derived functor of the Serre quotient embedding is an equivalence}
\author{The Stacks Project Authors}
\date{stacks-project:derived.tex:local-snapshot-20261012; GNU FDL 1.2 or later}
\begin{document}
\maketitle
\noindent\textit{Editorial note.} Counted span is the author's own proof environment, source lines 6370-6416, following the statement at 6360-6368.\par\medskip
\noindent\textit{Editorial note.} 'Fully faithfulness is proved in the exact same manner' (6414-6415): mirror of the executed faithfulness argument with the homotopy applied to the other composite, recorded as a routine repeat.\par\medskip
\noindent\textit{Allowed standard prerequisites (editorial).}\begin{itemize}
\item Serre subcategory closure. the Serre hypothesis of the statement, used at proof line 6399. Supplies that (D\textasciicircum{}\{i+1\} + E\textasciicircum{}\{i+1\}) intersect Im(d\textasciicircum{}i) lies in B, making the inductive choice of E\textasciicircum{}i possible.
\end{itemize}\medskip
\ifdefined\hypertarget\hypertarget{sourcepointer-1}{}\fi
\noindent\textit{Source reference.} Original equation 06UR: sources/stacks-derived/pinned-original-source.tex, source lines 6188--6188. Navigation pointer only: the counted proof cites this statement by \textbackslash{}ref as a standing input; the deduction the counted proof has to make is carried out inside the counted span, and no part of that deduction is delegated to this pointer.\par\medskip
\begin{lemma}
\label{lemma-fully-faithful-embedding}
Let $\mathcal{A}$ be an abelian category. Let $\mathcal{B} \subset \mathcal{A}$
be a Serre subcategory. Assume that for every surjection $X \to Y$
with $X \in \Ob(\mathcal{A})$ and $Y \in \Ob(\mathcal{B})$ there exists
$X' \subset X$, $X' \in \Ob(\mathcal{B})$ which surjects onto $Y$.
Then the functor $D^-(\mathcal{B}) \to D^-_\mathcal{B}(\mathcal{A})$ of
(\ref{equation-compare}) is an equivalence.
\end{lemma}

\begin{proof}
Let $X^\bullet$ be a bounded above complex of $\mathcal{A}$ such that
$H^i(X^\bullet) \in \Ob(\mathcal{B})$ for all $i \in \mathbf{Z}$.
Moreover, suppose we are given $B^i \subset X^i$, $B^i \in \Ob(\mathcal{B})$
for all $i \in \mathbf{Z}$. Claim: there exists a subcomplex
$Y^\bullet \subset X^\bullet$ such that
\begin{enumerate}
\item $Y^\bullet \to X^\bullet$ is a quasi-isomorphism,
\item $Y^i \in \Ob(\mathcal{B})$ for all $i \in \mathbf{Z}$, and
\item $B^i \subset Y^i$ for all $i \in \mathbf{Z}$.
\end{enumerate}
To prove the claim, using the assumption of the lemma we first choose
$C^i \subset \Ker(d^i : X^i \to X^{i + 1})$, $C^i \in \Ob(\mathcal{B})$
surjecting onto $H^i(X^\bullet)$. Setting
$D^i = C^i + d^{i - 1}(B^{i - 1}) + B^i$ we find a subcomplex
$D^\bullet$ satisfying (2) and (3) such that
$H^i(D^\bullet) \to H^i(X^\bullet)$ is surjective for all $i \in \mathbf{Z}$.
For any choice of $E^i \subset X^i$ with $E^i \in \Ob(\mathcal{B})$ and
$d^i(E^i) \subset D^{i + 1} + E^{i + 1}$ we see that setting
$Y^i = D^i + E^i$ gives a subcomplex whose terms are in $\mathcal{B}$ and
whose cohomology surjects onto the cohomology of $X^\bullet$. Clearly, if
$d^i(E^i) = (D^{i + 1} + E^{i + 1}) \cap \Im(d^i)$ then we see that
the map on cohomology is also injective. For $n \gg 0$ we can
take $E^n$ equal to $0$. By descending induction
we can choose $E^i$ for all $i$ with the desired property.
Namely, given $E^{i + 1}, E^{i + 2}, \ldots$ we choose $E^i \subset X^i$
such that $d^i(E^i) = (D^{i + 1} + E^{i + 1}) \cap \Im(d^i)$.
This is possible by our assumption in the lemma combined with
the fact that $(D^{i + 1} + E^{i + 1}) \cap \Im(d^i)$ is
in $\mathcal{B}$ as $\mathcal{B}$ is a Serre subcategory of $\mathcal{A}$.

\medskip\noindent
The claim above implies the lemma. Essential surjectivity is immediate
from the claim. Let us prove faithfulness. Namely, suppose we have
a morphism $f : U^\bullet \to V^\bullet$ of bounded above complexes
of $\mathcal{B}$ whose image in $D(\mathcal{A})$ is zero. Then
there exists a quasi-isomorphism $s : V^\bullet \to X^\bullet$
into a bounded above complex of $\mathcal{A}$ such that
$s \circ f$ is homotopic to zero. Choose a homotopy
$h^i : U^i \to X^{i - 1}$ between $0$ and $s \circ f$.
Apply the claim with $B^i = h^{i + 1}(U^{i + 1}) + s^i(V^i)$.
The resulting map $s' : V^\bullet \to Y^\bullet$
is a quasi-isomorphism as well and $s' \circ f$ is homotopic
to zero as is clear from the fact that $h^i$ factors through $Y^{i - 1}$.
This proves faithfulness. Fully faithfulness is proved in the
exact same manner.
\end{proof}

\end{document}
